\documentclass[11pt]{amsart}
\usepackage{amsmath,amssymb}
\usepackage[margin=1.1in]{geometry}
\usepackage{hyperref}
\title{Corrigenda and addenda to the November 2026 Gallai manuscripts}
\author{Compiled from independent referee audits}
\date{26 November 2026 (living document)}
\begin{document}
\begin{abstract}
Adversarial audits of the load-bearing parts of the collection found \emph{no false step} in the
reduction chain from Gallai's conjecture to the three-connected deletion tests, nor in the
block-degree-seven theorem, the first-failure alternative, or the odd semi-clique theorem, apart
from the external dependency recorded in Section~\ref{sec:jm}. This document lists every correction
found: wording and quantifier fixes, citation corrections, and one-line arguments that should be
written out. Each entry gives the location, the issue and the replacement text. The detailed audit
reports (\texttt{research/audit/A1}--\texttt{A5}) contain the verifications and computational checks.
\end{abstract}
\maketitle

\section{Cut vertices and cycle intersections (marked-reduction)}

\begin{enumerate}
\item \textbf{p.\,2, last sentence of Theorem~1.} ``In the first assertion it suffices to test nonempty
sets consisting entirely of even-degree vertices'' is proved (\S2.1) only \emph{in conjunction with}
assertion~2. Replace by: ``In the presence of assertion~2, the first assertion may be restricted to
nonempty sets of even-degree vertices.''
\item \textbf{p.\,2.} ``assembled from three or more parts'' should read ``from two or more parts'',
matching Theorem~7.
\item \textbf{p.\,3, lower bound in (2).} Give the reason as: $z$ is a cut vertex of $W(F,S)$, so a
path meets at most one copy of $X$ beyond $z$ and visits $z$ at most once.
\item \textbf{pp.\,2, 7, 12.} Define $q(F,\emptyset)$ (the path number of $F$) where the marked bound is
introduced; Lemma~6 uses it before \S2.1.
\item \textbf{p.\,4.} Insert the one-line proof of $\alpha_i=q(J_i^+,S_i\cup\{\ell_i\})$: an assignment
to the new leaf is the same as extending one path by the pendant edge.
\item \textbf{p.\,13, Theorem~7.} State that the induction uses the marked bound for \emph{all}
markings of smaller blocks, which is available by \S2.1.
\item \textbf{p.\,82, claim in Theorem~57.} Expand the integrality step: the two half-counts are
integers of the same parity, so the strict inequality improves by one.
\item \textbf{p.\,84, Theorem~59 finite-order clause.} State the conclusion explicitly: ``If both tests
hold for all members of $\mathcal C$ of order at most $m+1$, then the marked bound (1) holds on every
connected graph whose $2$-connected blocks have order at most $m$.'' Theorem~113 uses the case
$m=n-2$. The order-$(m+1)$ graphs involved are the cones $B+z$ and $B+w$.
\item \textbf{p.\,85, diamond lower bound.} Add: $d_B(x)\ge3$, which holds since $B$ is $2$-connected and
$d_B(x)$ is odd.
\item \textbf{p.\,105, Theorem~83, ``$m=3$'' exceptional bridge.} The case is impossible: if $x\neq r$
(here $r$ is the retained port, not the root) the retained side has a proper $2$-edge cut. The
construction that follows may be deleted. Rename the retained port to avoid the clash with the root.
\item \textbf{p.\,129, Theorem~114 proof.} Add the sentence covering the deletion of at most one vertex:
$B-v$ is connected and $z$ keeps a neighbour since $|S|\ge2$; and $J-z=B$.
\item \textbf{p.\,129, deletion-degree-six remark.} The consequence ``a failure at $d_J(z)=6$ needs two
universal vertices in $K$'' rests on Shaik's unrefereed preprint [25, Cor.~1.2]. Mark it as
conditional on [25], or replace the citation by an argument from Botler--Sambinelli, which is
refereed. Fan--Hou--Zhou alone gives one universal vertex.
\item \textbf{p.\,8 and p.\,129, Fan--Hou--Zhou [15].} Quote the exact published hypothesis. In
particular, say whether ``maximum degree $\le3$'' is required within each block or for the whole
$E$-subgraph; the doubling step of \S2.1 is sensitive to this.
\item \textbf{Fan's lemmas [14, Lemmas 3.3--3.4].} Quote them once verbatim, including the clause
``all other end counts are unchanged'', which every application uses.
\item \textbf{p.\,130.} ``$K_2\vee K_3$'' should read $K_2\vee\overline{K_3}$ (three independent pages).
\item \textbf{Theorem~144 can be made independent of [25].} At the step ``$r\in T$'' replace
[25, Cor.~1.2] by Botler--Sambinelli [7, Thm~13]: $B$ has four even vertices, so $\Delta(EV(B))\le3$, and
$B$ has even order, so it is not a SET graph. Replace ``By Corollary~106 with mark $r$, at least two
vertices of $K$ have degree four'' by ``by [7, Thm~13] together with Lemma~133''. Cite [7] in the two
remaining uses of Corollary~106 that have at most four unmarked even vertices. Theorem~144 then rests
only on refereed results. Corollary~106 itself, in its one-exception odd-order case, still uses [25].
\item \textbf{Lemma~133.} Its proof uses only connectivity of $B$. Weaken the hypothesis from
``two-connected'' to ``connected'', which is needed for the use above.
\item \textbf{Corollary~138} rests on the tables of Theorem~135 (9{,}330 entries) and Theorem~137 (235
cases) and on Lemma~121. Include them; they are now the main unreplayed inputs of the deletion-degree-six
programme.
\item \textbf{Theorem~124 (p.\,141).} The count ``10{,}906 graphs'' depends on the choice of labelled
representatives and of the route $L$, which the text does not fix. An independent replay obtained
13{,}482 under its own choices, while every invariant count (1{,}044 classes, 271 structural,
15{,}792 choices $(p,q,A)$, all with $|D|\le5$) agrees. Present 10{,}906 as implementation-dependent,
and publish the certificate files. Two structural seven-site shapes (a diamond, or $K_4$ joined by a bridge
to a triangle) use Fan--Hou--Zhou, which should be cited at that point.
\item \textbf{Lemma~126 (p.\,142).} Quote the Chu--Fan--Zhou star lemma [8, Lemma~2.1] that it uses.
\item \textbf{Addendum (Theorem~113).} The proof uses only that $B$ is $2$-connected and that both tests
hold on smaller members of $\mathcal C$. Stating it in this generality makes the later uses
(Theorems~144, 145, Corollary~146) transparent.
\item \textbf{Supplement.} The paper names \texttt{ApexConnectivity.lean}, \texttt{FiveSiteFan.lean},
\texttt{classify\_eight.cpp}, \texttt{check\_thin\_pairs.py} and others, but gives no location. Add a
repository reference. Independent audits re-derived the apex equivalence by hand and checked it by
computer through order~$8$.
\end{enumerate}

\section{Separators and endpoint exchange (separating-pairs)}

\begin{enumerate}
\item \textbf{pp.\,60--61, Theorem~12.6.} The ``unused slot'' is attributed to [6], a paper on
\emph{transversal} matroids, which does not cover the graphic matroid used here. The argument in the
text is correct; replace the citation by the direct argument or by a reference on graphic matroids.
\item \textbf{pp.\,77--78, Theorem~12.29.} The cycle-merging step cites Chu--Fan--Zhou [3, Lemma~2.4]
through a paraphrase of its hypotheses. The paper's own Lemma~13.6 covers the configuration; cite it
instead.
\item \textbf{Theorem~10.16.} The Eulerian clause ($\delta\ge6$) follows from $\delta\ge5$ and even
degrees; no separate argument is needed. Theorem~10.15 is not used in the proof of Theorem~10.16; it
enters only via Corollary~10.17 and later. Say so to clarify the dependency.
\item \textbf{Throughout.} Several results are cited by descriptive names (e.g.\ ``cubic odd-core
proposition''). Replace them by theorem numbers.
\item \textbf{References.} [15] (Lov\'asz 1968) and [16] (Hilton) are never cited in the text.
\end{enumerate}

\section{Odd semi-cliques and sharp path partitions (odd-semi-clique-boundary)}

\begin{enumerate}
\item \textbf{pp.\,9, 12: citation.} ``Javadi--Miralaei, Theorem~3.8'' does not exist in
arXiv:2403.17290v1; the matching-extension result is \textbf{Theorem~8}. The proof of Theorem~11 uses
Theorem~8 through its proof, which needs $r\le 2t-1$ rather than the stated $G\not\cong tK_2$. Our
setting satisfies $r\le2t-1$; say so explicitly.
\item \textbf{p.\,4: dependency bookkeeping for Theorem~4.} The statement that Javadi--Miralaei is needed
only for $k\ge6$, $\lceil (k+4)/2\rceil\le m\le k-1$ concerns the case $m<k$. The case $m=k$ (every
minimum partition Hamiltonian) needs it for every $k$-edge $H$ with $k\ge6$, $|V(H)|\ge k+2$ that is
not a linear forest.
\item \textbf{Notation $e(G)$.} In Theorems~20, 28, 31, Corollaries~23, 27, 33, 34, 40 and Proposition~26,
$e(G)$ denotes the number of \emph{missing} edges. Elsewhere it denotes $|E(G)|$. Replace the former by
$e(\overline G)$. Read literally, Corollary~27 is false ($K_{2k+1}$ is a counterexample), and
Theorem~20, Corollary~23 and Corollary~40 are vacuous. The proofs use the intended meanings correctly.
In Proposition~26 and Theorem~31, the set $X=\{v:d_G(v)>.01n\}$ should use the missing degree
$d_{\overline G}(v)$.
\item \textbf{p.\,45, Corollary~44.} The equivalence holds as a conjunction over all $C>0$. The hosts
used have about $|V(X)|\cdot N$ missing edges, so no single fixed $C$ suffices. Rephrase as: ``Gallai's
conjecture holds if and only if, for every $C>0$, it holds for all sufficiently large connected graphs
of even order $N$ with at most $CN$ missing edges.''
\end{enumerate}

\section{Blocks of maximum degree nine or cutwidth eight (block-degree)}

\begin{enumerate}
\item \textbf{Readiness statement.} Theorem~3 (block degree $\le7$) and Theorem~2 do not use
\S\S10--13. An independent replay reproduced every enumeration they use: the counts $494$; $63$; $1{,}348+989$;
$1{,}320$, $137$ and $03476125$ with $27$ sets; and the \S8 tables. Degree $8$, cutwidth $8$ and the
families of Theorems~7--13 depend on the \S10.2 computations (25{,}003 and 1{,}154 cases). Degree $9$
additionally depends on the \S12 and \S13 computations. The collection index should be corrected to say
this.
\item \textbf{p.\,4.} The five-contact path--cycle lemma is Zhang--Liu--Hong, Discrete Math.\ 2024
([17]). The two-cycle five-contact lemma is Chu--Fan--Liu. Both were re-proved exhaustively for at most
five contacts in the audit, so Theorem~3 no longer depends on these citations.
\item \textbf{p.\,4.} ``a path of arbitrary length'' should read ``of length at least $2$''.
\item \textbf{p.\,9, \S8.} State its hypothesis of block degree at most $7$; the incidence argument fails
at degree $8$.
\item \textbf{pp.\,5--7, 10.} ``Two different segments'' must allow both members of one parallel class.
\item \textbf{p.\,53.} The appendix check of Theorem~2 for seven-vertex components omits cases of maximum
degree above $7$ (paths through a cut vertex). It is not load-bearing.
\end{enumerate}

\section{Deletion-degree-six studies}

An independent exhaustive search re-derived the certificate banks of six-even-root-shapes (contact table
K0--K3 with 31/30/27/15 certified; all seven one-universal cores $\times$ 31 contact sets; all 32 masks of
the K2 refinement) and matched the printed survivors exactly. The banks should be published or replaced
by that search. Fan's Lemma~3.3 (grouped transfer) has been checked verbatim against the primary source,
and both of Fan's lemmas are now formally verified in Lean.

\begin{enumerate}
\item \textbf{six-even-incidence.} Cite Botler--Sambinelli (Thm 13) for the subcubic case. Cross-reference
six-clique Thm~1 for the circulant family, which is already covered there. Supply or remove the
``separate argument'' for separating clique edges.
\item \textbf{six-clique, Lemma~2.} Expand the chain step (Fan's backward uniqueness); also expand the
triple rectangles of Cor~5.
\item \textbf{degree-ten-marked-root.} Quote Chu--Fan--Zhou Lemma~2.4 exactly or replace it by a direct
finite check; include the selection-lemma and five-edge tables.
\end{enumerate}


\section{External dependency: Chu--Fan--Zhou Lemma~2.4}
The manuscripts use [Chu--Fan--Zhou, Lemma~2.4] in the form: a path $P$ with $|V(P)|\le4$ and $q$
edge-disjoint cycles meeting it, with $|V(P)|+q\le7$, partition into $q+1$ paths. The only verifiable
published statement (Chu--Wang, Lemma~2.2) has $|V(P)|+q\le6$. The note \texttt{manuscripts/cfz-lemma} proves
the needed form when every cycle passes through a common vertex of $P$, with $|V(P)|+q\le7$. It
contains one finite ingredient, a six-contact path--cycle lemma (1{,}257 compressed cases, agreeing with
\S6 of marked-reduction). Every use with sum seven satisfies the common-vertex hypothesis:
marked-reduction Lemma~21, Prop.~33, Cor.~35, Thm.~177, and even-degrees with $m=6$. The remaining uses need only
the sum-six form. At $m=7$, even-degrees should cite marked-reduction Theorem~38 instead. Cite the note at
these places.

\section{External dependency: Shaik's preprint (marked-reduction Corollary~106 and Theorem~144)}\label{sec:shaik}

Marked-reduction cites the unrefereed preprint of Shaik [hexagon:2610.00286v1] in four places:
\begin{itemize}
\item twice in the proof of Corollary~106;
\item in the remark after Theorem~114 (``at $d_J(z)=6$ a failure requires two universal vertices'');
\item at the step ``$r\in T$'' of Theorem~144.
\end{itemize}
Theorem~144 also inherits the preprint through Corollary~106. The note \emph{bs-prescribed}
(in \texttt{manuscripts/}) removes every one of these uses. Shaik's preprint is not needed anywhere in marked-reduction.

\begin{enumerate}
\item \textbf{Corollary~106, statement (p.\,120).} Add the hypothesis ``and the unmarked even sites number
at most five, or induce a graph of maximum degree at most three''. Every application in the manuscript
satisfies it: Theorem~144 three times (five, four and at most four unmarked even sites), and the
paragraph after Corollary~106 (seven even sites, two of them marked). The one-exception odd-order case
with at least seven unmarked even sites is not used, and it remains open without the preprint.
\item \textbf{Corollary~106, proof, first paragraph.} Replace the doubling argument by the following.
\begin{itemize}
\item If $|L|$ is even, ceiling and floor coincide. Use Botler--Sambinelli Cor.~14 (no exceptional
vertex) or Fan--Hou--Zhou Thm.~5 (one exceptional vertex).
\item If $|L|$ is odd, use Botler--Sambinelli Thm.~13 (no exceptional vertex; the leaf excludes SET
graphs) or bs-prescribed Cor.~2 (five even vertices, one of degree four).
\end{itemize}
The latter rests only on Lov\'asz's theorem and Fan's Lemmas~3.3--3.4. Six of the seven possible
even subgraphs are settled by three local lemmas. The bowtie is settled by eleven transfer certificates,
using a leaf at an odd support. All $131$ labelled certificates are checked in Lean
(\texttt{FiveEven.lean}, \texttt{FiveEvenCerts.lean}).
\item \textbf{Corollary~106, proof, case $A=\emptyset$.} Replace ``Shaik's ordinary bound'' by
``Botler--Sambinelli Cor.~14 or Fan--Hou--Zhou Thm.~5''.
\item \textbf{Remark after Theorem~114 (p.\,129).} Replace ``Shaik's ordinary bound [25, Cor.~1.2] also
applies \ldots\ since doubling gives at most two'' by: ``At $d_J(z)=6$ both tests are floor bounds for
$J-z$ with one or two private leaves at $r$. These graphs have odd order, leaves at the odd support $r$,
and the five even vertices $N(z)\setminus\{r\}$. So bs-prescribed Cor.~2 applies when $K$ has at most one
universal vertex.''
\item \textbf{Theorem~144, step $r\in T$ (p.\,156).} Replace [25, Cor.~1.2] by Botler--Sambinelli Thm.~13.
Alternatively use endpoint-structure Cor.~22, which needs only Lov\'asz and Fan: $B$ has four even
sites and $r$ is odd.
\item \textbf{Theorem~144, ``By Corollary~106 with mark $r$''.} This is now bs-prescribed Cor.~2 (five
even sites). The alternative of audit A6 (Lemma~133 together with its absent M/P/C table) is no longer
needed.
\end{enumerate}

\emph{Why not a prescribed-end version of Botler--Sambinelli?} That route was proposed in the project
notes (connections~\S L), and it is false. There are counterexamples at every odd order $\ge5$, including
graphs where the prescribed vertex has degree two between two odd vertices and the even subgraph is a
matching. The joining gadget also fails on an order-nine instance of the very case needed. See
bs-prescribed \S3.

\section{External dependency: the rainbow Hamiltonian decomposition theorem}\label{sec:jm}

The odd semi-clique theorem (Theorem~4 of odd-semi-clique-boundary and the September note), its
endpoint refinement (Theorem~11), the equality case inside Theorem~10.16 of separating-pairs, and
several other results use the theorem of Javadi and Miralaei [Discrete Math.\ 349 (2026) 114676]: every
$k$-edge subgraph of $K_{2k+1}$ is rainbow in some Hamiltonian decomposition. The audit found that its
proof (Theorem~7 combined with Theorem~8) has
\begin{itemize}
\item a fixable gap in the transfer steps of Theorem~7. The fix is to orient the paths and delete
out-edges at receiver ends; the counts are unchanged;
\item a substantive gap in Claim~12 of Theorem~8. The pairing step can fail when both duplicates of an
endpoint edge survive, and condition (iv) does not imply the linear-forest conclusion.
\end{itemize}
No counterexample is known. The theorem was verified by computer for every $k$-edge subgraph up to
isomorphism with $k\le7$. Liu and Chen proved it when the subgraph has at most $k+1$ vertices, is a
linear forest, or $k\le5$. \textbf{Until a repaired proof is confirmed, every result using the theorem
outside these ranges should state the dependency explicitly.} The three-connected reduction
(Theorems~1--114 of marked-reduction) and the block-degree-seven theorem do not depend on it. A repair
is in preparation.

\end{document}
